\usepackage{iftex}
\ifPDFTeX
  \usepackage[utf8]{inputenc}
  \usepackage[T5]{fontenc}
  \usepackage{lmodern}
\else
  \usepackage{fontspec}
\fi
\usepackage[vietnamese]{babel}
\usepackage[top=2cm,bottom=2cm,left=2cm,right=2cm]{geometry}
\usepackage{amsmath,amssymb}
\usepackage{graphicx}
\usepackage{tikz}
\usepackage{float}
\usepackage{caption}
\usepackage{enumitem}
\usepackage{fancyhdr}
\usepackage{lastpage}
\usepackage{booktabs}
\usepackage{colortbl}
\usepackage{array}
\usepackage[colorlinks=true,urlcolor=blue,linkcolor=black]{hyperref}
\usepackage[most]{tcolorbox}
\usepackage{needspace}

\newtcolorbox{dapsobox}[1][]{
  colback=blue!3!white,
  colframe=blue!60!black,
  fonttitle=\bfseries,
  title=Đáp số,
  arc=1.5mm,
  boxrule=0.8pt,
  left=8pt,
  right=8pt,
  top=6pt,
  bottom=6pt,
  #1
}

\captionsetup{font=it,labelfont=it}

\pagestyle{fancy}
\fancyhf{}
\fancyfoot[L]{\itshape \courseName}
\fancyfoot[R]{\itshape Trang \thepage}
\renewcommand{\headrulewidth}{0pt}
\renewcommand{\footrulewidth}{0.4pt}

\renewcommand{\thesection}{\arabic{section}}
\setcounter{secnumdepth}{2}

\newcommand{\notes}{}

% Ô dán ảnh: tự hiện ảnh nếu tệp tồn tại, nếu chưa có thì hiện khung màu cam để biết chỗ dán
\newcommand{\figslot}[4]{%
  \begin{figure}[H]
    \centering
    \IfFileExists{#1}{\includegraphics[#2]{#1}}{%
      \begin{tcolorbox}[enhanced,center,width=0.92\linewidth,height=#4,valign=center,halign=center,
          colback=yellow!8,colframe=orange!80!black,boxrule=1pt,arc=1.5mm]
        \textbf{Ô DÁN ẢNH}\\[2pt]
        {\small\itshape Lưu ảnh thành tệp\\ \texttt{\detokenize{#1}}\\ (khung này tự biến mất khi có ảnh)}
      \end{tcolorbox}}%
    \caption{#3}
  \end{figure}}

\newcommand{\given}[3]{\figslot{images/given/#1.png}{width=#2,height=0.32\textheight,keepaspectratio}{#3}{5cm}}
% \simfig[chiều cao ô]{tên tệp trong images/sim}{chú thích}
\newcommand{\simfig}[3][6.5cm]{\figslot{images/sim/#2.png}{width=0.94\linewidth,height=0.45\textheight,keepaspectratio}{#3}{#1}}
% ảnh sơ đồ mạch có hiển thị V, I (khổ nhỏ hơn ảnh Probe)
\newcommand{\simfigs}[2]{\figslot{images/sim/#1.png}{width=0.72\linewidth,height=0.28\textheight,keepaspectratio}{#2}{5cm}}
\newcommand{\drawfig}[2]{\figslot{images/draw/#1.png}{width=0.9\linewidth,height=0.45\textheight,keepaspectratio}{#2}{4cm}}

% Ô số liệu mô phỏng trong bảng: thay \fillcell bằng giá trị đọc được
\newcommand{\fillcell}{\fcolorbox{orange!80!black}{yellow!25}{\rule{0pt}{1.3em}\rule{1.05cm}{0pt}}}

% Ô ghi nhận xét / quan sát của sinh viên
\newtcolorbox{notebox}[2][]{enhanced,colback=gray!7,colframe=gray!65!black,fonttitle=\bfseries,
  title={#2},boxrule=0.8pt,arc=1.5mm,left=8pt,right=8pt,top=6pt,bottom=6pt,#1}

\newcommand{\simresult}{\Needspace{0.55\textheight}\par\noindent\textbf{\textit{Kết quả mô phỏng (hình ảnh):}}\par}

\newcommand{\dotline}[1]{\noindent #1\dotfill\par\vspace{0.8em}}
\newcommand{\dotlineeq}[2]{\noindent #1\dotfill\ #2\par\vspace{0.8em}}
\newcommand{\dotlines}[1]{\foreach \i in {1,...,#1}{\noindent\dotfill\par\vspace{1em}}}
\newcommand{\elemline}[1]{\noindent #1\dotfill\ $\longrightarrow$\ \dotfill\ element\par\vspace{0.8em}}
